\section{Nguồn dữ liệu và pipeline xử lý}

\subsection{Nguồn dữ liệu}

Dữ liệu được nạp từ bốn tệp kết quả Excel:
\begin{itemize}[leftmargin=1.2cm]
    \item \texttt{results/random.xlsx}
    \item \texttt{results/reverse.xlsx}
    \item \texttt{results/nearly-sorted.xlsx}
    \item \texttt{results/many\_duplicates.xlsx}
\end{itemize}

Mỗi tệp chứa kết quả cho ba thuật toán trên các kích thước đầu vào:
\[
    n \in \{100, 1000, 10000, 100000\}
\]

\subsection{Schema chuẩn hóa}

Do dữ liệu gốc ở dạng bảng trình bày, pipeline chuyển về dạng tidy để phân tích và vẽ:
\begin{table}[H]
\centering
\renewcommand{\arraystretch}{1.2}
\begin{tabular}{lll}
\toprule
\textbf{Cột} & \textbf{Kiểu dữ liệu} & \textbf{Ý nghĩa} \\
\midrule
\texttt{input\_type} & category & Kiểu dữ liệu đầu vào \\
\texttt{algorithm} & category & Tên thuật toán \\
\texttt{n} & integer & Kích thước đầu vào \\
\texttt{time\_ms} & float & Thời gian chạy (ms) \\
\texttt{comparisons} & float/optional & Số phép so sánh (nếu có) \\
\bottomrule
\end{tabular}
\caption{Schema dữ liệu chuẩn hóa dùng cho trực quan hóa}
\label{tab:tidy-schema}
\end{table}

\subsection{Pipeline tạo figure}

Pipeline được triển khai trong mã nguồn \texttt{build\_figures.py} với các bước:
\begin{enumerate}[leftmargin=1.2cm]
    \item Đọc và parse các tệp Excel theo từng kiểu dữ liệu.
    \item Chuẩn hóa tên thuật toán, thứ tự category và nhãn biểu đồ.
    \item Sinh bảng \texttt{tidy\_results.csv} và \texttt{summary\_time\_ms.csv}.
    \item Vẽ và xuất figure dưới dạng raster (PNG) và vector (PDF, SVG).
\end{enumerate}

\subsection{Đảm bảo tính nhất quán}

Để so sánh công bằng giữa các biểu đồ, bài làm cố định các yếu tố sau:
\begin{itemize}[leftmargin=1.2cm]
    \item Cùng bảng màu, marker và thứ tự thuật toán trên mọi figure.
    \item Cùng mốc trục $x$ theo $n$ và cùng chú giải (legend) nhất quán.
    \item Cùng đơn vị đo thời gian là millisecond (ms).
\end{itemize}
